\begin{abstract}
Urban environmental risk research remains deeply fragmented into single-domain silos: meteorological frameworks monitor atmospheric extremes while treating the urban fabric as a featureless plane, whereas urban planning indices map static built infrastructure or social deprivation while treating climate variability as static or uniform. Crucially, neither paradigm captures how dynamic biophysical hazards non-linearly compound within dense, impervious urban morphologies. Here, we present the \textbf{Compound Sustainability Stress Index (SSI)}, a unified, spatially explicit multi-hazard framework that bridges continuous earth observation reanalysis with micro-scale urban built environment analytics. Grounded in a strictly balanced panel of \textbf{334,332 spatio-temporal hexagonal cells} (Uber H3 Resolution~8, $\sim 0.74\text{ km}^2$) monitored continuously across 36 months (2020--2022), the framework encompasses seven contrasting Indian metropolitan biomes (Bikaner, Coimbatore, Delhi, Hyderabad, Indore, Kolkata, and Mumbai). The model integrates five un-normalized physical drivers spanning physiological thermal sultriness (ERA5 Heat Index), hydrological deficit (Climatic Water Deficit), photosynthetic canopy depletion (Sentinel-2 NDVI), atmospheric particulate aerosols (Copernicus CAMS $\text{PM}_{2.5}$), and physical built exposure (GHSL impervious surfaces and OpenStreetMap building and road networks). To model multi-hazard synergy without mathematical compensation artifacts, we introduce a smooth asymptotic bivariate joint-tail hazard amplifier ($\gamma = 0.30$, $q = 0.90$) that captures a \textbf{+25.3\% compound stress escalation} during concurrent extremes while eliminating artificial boundary ceiling truncation ($0.00\%$ records clamped at unity). Unsupervised $k$-means clustering across the multi-biome panel reveals six stable pan-India urban stress archetypes, providing an actionable typology for municipal climate adaptation. An eight-part pre-registered validation battery verifies indicator orthogonality ($\text{VIF} < 1.12$), cross-biome differentiation ($F = 6,888$, $p < 10^{-15}$), and precise temporal alignment with recorded historical disasters, including the Spring 2022 Northwest India mega-heatwave (96.5\% alert rate in Bikaner) and the October 2021 catastrophic Hyderabad flash floods (100\% alert rate). Furthermore, spatial prioritization exhibits over $94\%$ rank invariance across alternative weighting regimes (Equal, Entropy, CRITIC, and PCA). By breaking single-domain silos, this framework delivers an open, scalable diagnostic for compound climate resilience across rapidly urbanizing landscapes in the Global South.
\end{abstract}

\noindent\textbf{Keywords:} Compound Climate Hazards; Urban Sustainability; Spatio-Temporal Modeling; Discrete Global Grid Systems; Multi-Hazard Synergy; Global South.
